\chapter{Wnioski}

\section{Podsumowanie eksperymentów}

Zakres eksperymentu objął dostrajanie hiperparametrów pięciu wariantów modeli implicit ALS z różnymi wagami pewności (liniowa, logarytmiczna, logarytmiczna z saturacją, logarytmiczna z przesunięciem $\beta$ oraz stała), modelu BPR oraz popularnościowego jako odniesienie. Dobór hiperparametrów pokazał, że w przypadku modeli iALS optymalny hiperparametr $\alpha \to 0$, co w praktyce redukowało sygnał o preferencji płynący z \texttt{playtime\_forever} do wariantu binarnego (użytkownik posiada/nie posiada gry w bibliotece). Jedynym wyjątkiem był model z $c_{ui} = 1 + \alpha \log(1 + r_{ui} + \beta)$, w którym optymalna wartość $\alpha$ nie leżała w otoczeniu zera. Analiza kształtu tego przekształcenia wykazała jednak, że przy $\beta$ równym 108,284 funkcja działa w obszarze o niewielkiej krzywiźnie: dla mediany czasu gry wynoszącej 35 minut wartość $c_{ui}$ wynosi 2,94 wobec 2,83 dla obserwacji bez zarejestrowanego czasu rozgrywki. Wyższa wartość $\alpha$ podnosiła zatem przede wszystkim ogólny poziom pewności, w mniejszym stopniu różnicując obserwacje między sobą.

Obserwacja ta doprowadziła do zbadania granicznego przypadku całej rodziny, czyli wariantu $c_{ui} = 1 + \gamma$, całkowicie pomijającego czas gry. Wariant ten osiągnął najwyższe wyniki metryk rankingowych zarówno na zbiorze walidacyjnym, jak i testowym, poprawiając \ndcg{} o 6,2\% względem wariantu logarytmicznego oraz o 8,1\% względem liniowego. Osiągnął on również największe pokrycie katalogu wśród modeli implicit ALS. Wszystkie różnice okazały się istotne statystycznie (test Wilcoxona z korektą Holma-Bonferroniego, poziom istotności $= 0{,}05$), przy czym analiza wielkości efektu wykazała, że mediana sparowanej różnicy \ndcg{} we wszystkich porównaniach wynosi zero. Obserwowana przewaga ma zatem charakter zagregowany i nie przekłada się na przewidywalną poprawę jakości rekomendacji dla pojedynczego użytkownika.

Badanie pokazało również kompromisy między jakością rankingową a miarami wykraczającymi poza trafność. Model BPR uzyskał ponad dwukrotnie większe pokrycie katalogu niż najlepszy z modeli iALS oraz wysoką różnorodność wewnątrzlistową, ale osiągnął to kosztem znacznie niższej trafności rekomendacji. Natomiast wariant liniowy, mimo najniższej jakości rankingowej w rodzinie iALS, cechuje się największą różnorodnością rekomendacji oraz zdolnością polecania obiektów mniej popularnych niż średnia popularność w zbiorze treningowym danego użytkownika, co jest nietypowym zachowaniem dla modelu faktoryzacji macierzy.

Przeprowadzona analiza stabilności ze względu na losowe uruchomienia algorytmów oraz podziały danych wykazała wysoką powtarzalność wyników. Odchylenia standardowe metryk trafnościnie przekraczały 0,2\% wartości średniej, co oznacza, że obserwowane różnice między wariantami nie wynikają z~losowości procedury.

Powyższe wyniki pozwalają odpowiedzieć na postawione pytania badawcze.

\section{Odpowiedź na pytania badawcze}

\textbf{PB1. Czy sposób przekształcenia czasu gry na wagi pewności wpływa na jakość rekomendacji modelu implicit ALS?}

Tak, jednak w~kierunku odwrotnym do oczekiwanego. Różnica \ndcg{} między najlepszym a~najgorszym wariantem wynosi 0,0266, co odpowiada poprawie o~8,1\% względem wariantu liniowego. Analiza sparowanych różnic na poziomie użytkownika potwierdza istotność tego efektu: średnia różnica wynosi 0,0266 przy 95-procentowym przedziale ufności $[0{,}0248;\ 0{,}0284]$. Wyższe wyniki osiągnęły jednak te przekształcenia, które w mniejszym stopniu uzależniały wagi pewności od zarejestrowanego czasu rozgrywki. Najwyższą jakość rekomendacji uzyskał wariant całkowicie pomijający tę zmienną, najniższą zaś wariant liniowy, przypisujący jej wagi proporcjonalne do czasu gry. Należy przy tym podkreślić, że porównywane warianty dostrojono niezależnie, przez co różnią się również liczbą czynników ukrytych oraz siłą regularyzacji. Wniosek dotyczy zatem tego, który wariant po dostrojeniu prowadzi do najlepszych rekomendacji, nie zaś izolowanego wpływu samego przekształcenia sygnału.

\textbf{PB2. Czy różnicowanie wag pewności według czasu gry przewyższa wariant o stałej wadze, niezależnej od intensywności interakcji?}

Nie. Wariant o stałej funkcji pewności uzyskał najwyższą wartość \ndcg{} na zbiorze testowym (0,3565), wyprzedzając wariant z przesunięciem $\beta$ (0,3517) oraz pozostałe przekształcenia. Dwa najlepsze warianty łączy podniesienie pewności obserwacji o zerowym czasie gry, którym pozostałe przekształcenia przypisują wartość równą 1. W badanym zbiorze monotoniczne uwzględnienie zmiennej \texttt{playtime\_forever} nie przyniosło zatem wartości dodanej wykraczającej poza sam wzrost wagi zaobserwowanych par użytkownik---gra.

\textbf{PB3. W jaki sposób wybór funkcji pewności przekłada się na miary wykraczające poza trafność i czy występuje kompromis między trafnością a tymi miarami?}

Odpowiedź różni się w zależności od rozważanej miary. W przypadku pokrycia katalogu kompromis w obrębie rodziny implicit ALS nie występuje: warianty o najwyższej trafności osiągnęły jednocześnie najwyższe pokrycie (0,1374 oraz 0,1313 wobec 0,0762 dla wariantów logarytmicznych). Wartości $\Delta$GAP wskazują natomiast, że wszystkie warianty poza liniowym rekomendują pozycje popularniejsze niż wynikałoby to z profilu użytkownika, przy czym dwa czołowe warianty wykazują umiarkowane obciążenie popularnością (0,4210 oraz 0,4191) wobec wartości powyżej 0,54 dla wariantów logarytmicznych. Warto podkreślić, że ujemna wartość $\Delta$GAP wariantu liniowego ($-0{,}2112$) stanowi odchylenie od punktu neutralnego tak samo jak wartości dodatnie, jedynie w przeciwnym kierunku, i również nie jest pożądana. W~przypadku różnorodności wewnątrzlistowej wartości uzyskane w~rodzinie implicit ALS układają się zgodnie z~liczbą czynników ukrytych, dobieraną niezależnie dla każdego wariantu. Nie jest to ograniczenie samej metryki, która wyznaczana jest we~wspólnej dla wszystkich modeli przestrzeni współwystępowania, lecz obserwacja dotycząca zachowania modeli --- model BPR mimo najwyższej liczby czynników osiąga różnorodność niższą niż wariant liniowy. Rozdzielenie wpływu wymiarowości od wpływu samej funkcji pewności wymagałoby jednak eksperymentu przy ustalonej liczbie czynników.

\textbf{PB4. Czy o jakości rekomendacji decyduje w większym stopniu sposób ważenia obserwacji, czy klasa modelu i postać funkcji celu?}

Uzyskane wyniki wskazują na przewagę drugiego z tych czynników. Różnica \ndcg{} między najlepszym wariantem implicit ALS a modelem BPR wynosi 0,0862 i jest ponad trzykrotnie większa niż cały rozstęp wewnątrz rodziny implicit ALS, obejmujący wartości od 0,3299 do 0,3565. Oba porównywane modele traktują przy tym interakcje w sposób nieuwzględniający czasu gry, a~mimo to dzieli je znacznie większa różnica niż którekolwiek dwa warianty funkcji pewności. Także i~w~tym przypadku porównanie dotyczy modeli dostrojonych niezależnie, obejmuje zatem łączny efekt zmiany funkcji celu i~towarzyszącego jej doboru pozostałych hiperparametrów. Odpowiedź nie jest jednak jednoznaczna: model BPR uzyskał pokrycie katalogu na poziomie 0,3039, ponad dwukrotnie przewyższające najlepszy wariant implicit ALS. Wybór klasy modelu determinuje zatem nie tyle poziom jakości rekomendacji, ile profil kompromisów między poszczególnymi jej wymiarami.

\section{Najważniejsze wnioski praktyczne i metodologiczne}

\subsection{Wnioski praktyczne}

Proste przekształcenia funcjkji pewności w przypadku danych o dużym zakresie są podatne na dwa przeciwstawne błędy. Liniowa funkcja jest wrażliwa na wartości odstające, natomiast przekształcenia logarytmiczne są wrażliwe na szum wśród małych wartości. W obu przypadkach proces optymalizacji hiperparametrów reagował w ten sam sposób, dobierając wartości $\alpha$ na tyle małe, że przekształcenie sprowadzało się w praktyce do postaci niemal stałej.

Sam fakt posiadania gry w bibliotece stanowi silny sygnał preferencji użytkownika. W badanym zbiorze okazał się on wystarczający do uzyskania najlepszych rekomendacji, a próby wydobycia dodatkowej informacji z długości interakcji nie przyniosły poprawy. Kluczowe znaczenie ma natomiast poziom wagi przypisanej zaobserwowanym interakcjom, stanowiący pełnoprawny hiperparametr wymagający strojenia niezależnie od postaci przekształcenia.

Wybór modelu powinien wynikać z~przyjętych celów systemu rekomendacyjnego. Jeżeli priorytetem jest trafność rekomendacji, wariant iALS o stałej funkcji pewności jest najlepszym wyborem, oferując jednocześnie najwyższe pokrycie katalogu w swojej rodzinie. W przypadku gdy cenniejsze jest szerokie pokrycie katalogu oraz zróżnicowanie rekomendacji między użytkownikami, warto rozważyć model BPR, akceptując wyraźnie niższą trafność.

Redukcja regularyzacji i zwiększenie wymiaru przestrzeni cech ukrytych $f$ może zmniejszać obciążenie popularnością $\Delta$GAP, co zaobserwowano w przypadku wariantu liniowego ALS oraz modelu BPR, osiągających ujemne wartości tej metryki. Należy jednak zachować ostrożność przy interpretacji tego efektu: wariant liniowy uzyskał jednocześnie najwyższą różnorodność rekomendacji i najniższą trafność w swojej rodzinie, co sugeruje, że obserwowane wartości mogą wynikać ze słabszego dopasowania modelu, a nie z lepszego rozpoznania niszowych preferencji.

\subsection{Wnioski metodologiczne}

Dokładna analiza wrażliwości hiperparametrów pozwoliła zinterpretować nie tylko to, co osiąga dany model, ale również dlaczego otrzymywane są takie wyniki. Dodatkowo wizualizacja funkcji pewności wraz z zakresem przyjmowanych przez nie wartości umożliwiła wykrycie, że przekształcenia formalnie różnicujące sygnał działają w praktyce w wąskim przedziale, czego sama wartość metryki nie ujawnia.

Szczególnie istotne okazało się uwzględnienie wariantu granicznego, pomijającego badaną zmienną w całości. Bez niego obserwowana poprawa kolejnych przekształceń zostałaby przypisana lepszemu odwzorowaniu siły preferencji, podczas gdy w rzeczywistości wynikała z podniesienia ogólnego poziomu pewności.

Nie mniej ważna okazała się konieczność wielowymiarowej oceny modeli rekomendacyjnych, uwzględniającej zarówno metryki trafności, jak i metryki wykraczające poza nią. Poleganie wyłącznie na \ndcg{} daje niepełny obraz jakości systemu rekomendacyjnego. Istotny jest przy tym dobór samych metryk: pierwotnie rozważana miara różnorodności, oparta na odległościach w przestrzeni czynników ukrytych, okazała się nieporównywalna między modelami o odmiennej wymiarowości, co wymusiło jej zastąpienie miarą opartą na współwystępowaniu obiektów we wspólnej dla wszystkich modeli przestrzeni.

Konieczne są również testy stabilności modelu ze względu na losowość podziału danych oraz inicjalizacji algorytmów, jak i testy istotności statystycznej różnic. Przy dużej liczbie użytkowników sama istotność statystyczna nie jest jednak wystarczającą podstawą do wnioskowania i wymaga uzupełnienia miarami wielkości efektu wraz z przedziałami ufności.

\section{Kierunki dalszych badań}

Kierunkiem wynikającym bezpośrednio z uzyskanych wyników, stanowiącym autorską propozycję, jest zbadanie niemonotonicznych funkcji pewności. Wszystkie przekształcenia rozważane w niniejszej pracy były monotonicznie rosnące, przyjmując założenie, że dłuższy czas gry oznacza silniejszą preferencję. Założenie to może jednak nie odpowiadać rzeczywistym zachowaniom użytkowników. W szczególności krótka interakcja, obejmująca kilkanaście czy kilkadziesiąt minut rozgrywki, może świadczyć nie o słabej preferencji, lecz o uruchomieniu gry i szybkiej rezygnacji z niej, stanowiąc tym samym sygnał negatywny. W takim ujęciu obserwacje o zerowym czasie gry, odpowiadające samemu posiadaniu tytułu w bibliotece, niosłyby silniejszy sygnał preferencji niż obserwacje o bardzo krótkim czasie rozgrywki. Funkcja pewności odzwierciedlająca ten wzorzec musiałaby przypisywać obniżone wagi w wąskim przedziale niskich wartości \texttt{playtime\_forever}, a rosnące dopiero powyżej pewnego progu. Zbadanie takiego wariantu pozwoliłoby zweryfikować, czy brak wartości dodanej z czasu gry zaobserwowany w niniejszej pracy wynika z samej zmiennej, czy z założenia o monotonicznym charakterze jej związku z preferencją.

Wyniki metryk $\Delta$GAP i pokrycia pokazały, że modele ALS (poza wariantem liniowym) mają tendencję do faworyzowania popularnych pozycji. Potencjalnym rozwiązaniem byłoby wprowadzenie do funkcji $c_{ui}$ lub do funkcji straty czynnika odwrotnie proporcjonalnego do popularności danego obiektu. Metody tego typu są opisywane w literaturze jako debiasing z wykorzystaniem odwrotnego ważenia prawdopodobieństwa ekspozycji (ang. \textit{inverse propensity scoring}, IPS), gdzie waga interakcji jest skalowana odwrotnie do prawdopodobieństwa, że dany obiekt zostałby zaobserwowany przez użytkownika niezależnie od jego rzeczywistych preferencji~\cite{schnabel2016recommendations}. W praktyce prawdopodobieństwo to jest zwykle estymowane na podstawie popularności obiektu, ponieważ użytkownicy z większym prawdopodobieństwem wchodzą w interakcję z obiektami popularnymi niezależnie od indywidualnych preferencji. Taka transformacja mogłaby podnosić wagę interakcji z niszowymi grami, kompensując mechanizm, przez który obecne funkcje marginalizują rzadkie obserwacje.

Poza modyfikacją samej funkcji pewności, alternatywnym kierunkiem jest wykorzystanie dopełniających się właściwości testowanych rodzin modeli poprzez ich połączenie w model hybrydowy. Przykładowo model BPR generowałby szeroki wachlarz kandydatów ze względu na większe pokrycie zbioru, a następnie przeprowadzany byłby ponowny ranking przez model iALS cechujący się większą trafnością.

Uzupełnieniem kierunków rozwoju wynikających z~głównych wniosków pracy są badania wskazane jako ograniczenia przyjętej metodologii:

\begin{itemize}
    \item segmentację użytkowników według liczby posiadanych interakcji,
    umożliwiającą zweryfikowanie hipotezy o~wyższej precyzji wariantu liniowego
    dla użytkowników o~większej aktywności;

    \item eksperyment przy ustalonej, jednakowej liczbie czynników ukrytych dla
    wszystkich wariantów, pozwalający rozdzielić wpływ wymiarowości od wpływu
    samej funkcji pewności na metryki różnorodności i~pokrycia;

    \item powtórzenie eksperymentu na danych ograniczonych do interakcji
    z~niezerowym czasem rozgrywki, stanowiących około dwóch trzecich zbioru,
    w~celu weryfikacji, czy brak wartości dodanej wynika z~właściwości samej
    zmiennej, czy ze sposobu jej wykorzystania;

    \item porównanie dostrojonego wariantu implicit ALS z~modelami neuronowymi,
    nawiązujące do wyników przedstawionych w~pracy~\cite{rendle2021revisiting};

    \item walidację uzyskanych wyników na innych zbiorach danych i~domenach,
    pozwalającą ustalić, czy zaobserwowany brak przewagi wariantów
    różnicujących jest właściwością uniwersalną, czy wynika z~charakterystyki
    rozkładu długości interakcji w~badanym zbiorze.
\end{itemize}